% Both url and microtype are loaded by acmart; pass options before it does.
\PassOptionsToPackage{hyphens}{url}
\PassOptionsToPackage{protrusion=true,expansion=false}{microtype}
\documentclass[sigconf,10pt]{acmart}

\usepackage{listings}
\usepackage{xcolor}
\usepackage{booktabs}
\usepackage{tabularx}

% Give TeX up to 3 em of extra stretch before declaring an overfull box.
% Combined with microtype's hz-optimization this eliminates almost all
% column-overflow without resorting to \sloppy's loose spacing.
\emergencystretch=3em

\lstset{
  basicstyle=\ttfamily\small,
  keywordstyle=\color{blue}\bfseries,
  commentstyle=\color{gray},
  breaklines=true,
  frame=single,
  numbers=left,
  numberstyle=\tiny\color{gray},
  xleftmargin=1.5em,
  framexleftmargin=1.5em,
}

\setcopyright{none}
% Suppress acmart's auto-filled conference/DOI placeholder block in the
% header and the copyright/permission footnote at the bottom of column 1.
\settopmatter{printacmref=false}
\renewcommand\footnotetextcopyrightpermission[1]{}
\AtBeginDocument{\fancyhead{}}

% ── license ───────────────────────────────────────────────────────────────────
% This work is licensed under Creative Commons Attribution 4.0 International
% (CC BY 4.0). See LICENSE-PAPER for the full legal text.

% ── metadata ────────────────────────────────────────────────────────────────
\title{openssl-provider-vault: Transparent Cryptographic Key Isolation\\
       for OpenSSL~3 via HashiCorp Vault}

\author{Soren L.~Hansen}
\email{sorenisanerd@gmail.com}
\affiliation{
  \institution{Independent Researcher}
  \city{}
  \country{}
}

\begin{abstract}
Private key exposure remains a primary cause of large-scale cryptographic
failures. Hardware Security Modules mitigate this risk but impose high cost and
operational complexity.  Software-based secret-management systems such as
HashiCorp Vault offer strong logical and network-level isolation guarantees
at a fraction of the deployment friction---when deployed on a separate
host, an adversary must compromise both the application process and the
Vault cluster simultaneously to extract key material, a strictly stronger
requirement than a single-system boundary---yet integrating them with
applications that rely on standard TLS and signature APIs has historically
required intrusive code changes.

We present \emph{openssl-provider-vault}, an OpenSSL~3 provider that routes
private-key operations such as signing and asymmetric encryption, together
with HMAC operations, through Vault's transit secrets engine without any
modification to applications or system libraries.  The provider implements
the standard OpenSSL~3 provider dispatch interface, making it a drop-in
provider for the supported key-management, signing, asymmetric-cipher, MAC,
and STORE operations.  We describe the security architecture, the design
trade-offs imposed by the OpenSSL~3 API, and lessons learned from bridging
the impedance mismatch between a streaming cipher API and a prehash-oriented
remote service.
\end{abstract}

\keywords{key management, OpenSSL provider, HashiCorp Vault, key isolation,
transit encryption, cryptographic API}

%%
%% CCS concepts
%%
\begin{CCSXML}
<ccs2012>
  <concept>
    <concept_id>10002978.10002986.10002990</concept_id>
    <concept_desc>Security and privacy~Key management</concept_desc>
    <concept_significance>500</concept_significance>
  </concept>
  <concept>
    <concept_id>10002978.10002986.10003017</concept_id>
    <concept_desc>Security and privacy~Cryptography</concept_desc>
    <concept_significance>300</concept_significance>
  </concept>
  <concept>
    <concept_id>10011007.10011006.10011039</concept_id>
    <concept_desc>Software and its engineering~Software libraries and repositories</concept_desc>
    <concept_significance>100</concept_significance>
  </concept>
</ccs2012>
\end{CCSXML}

\ccsdesc[500]{Security and privacy~Key management}
\ccsdesc[300]{Security and privacy~Cryptography}
\ccsdesc[100]{Software and its engineering~Software libraries and repositories}

% ────────────────────────────────────────────────────────────────────────────
\begin{document}
\maketitle

% ── 1. Introduction ──────────────────────────────────────────────────────────
\section{Introduction}
\label{sec:intro}

The compromise of a private key is a catastrophic, often irreversible event.
Once an adversary obtains the private key of a TLS certificate, a code-signing
identity, or a long-lived authentication credential, they can silently
impersonate the legitimate holder indefinitely.  Notable incidents---from
rogue certificate issuance to software supply-chain attacks---have
involved leaked or compromised private keys as a central element.

The classical defense is an \emph{Hardware Security Module} (HSM): a
tamper-resistant device that holds the key and exposes only high-level
operations (sign, decrypt).  The private key never leaves the HSM boundary
in cleartext.  The drawback is cost---enterprise HSMs range from several
hundred to several thousand dollars per unit---and operational complexity.
Deployments must manage physical devices, driver stacks, firmware updates,
and the coarse-grained concurrency model imposed by PKCS\#11.

Software-defined secret-management systems, and HashiCorp Vault in particular,
replicate the same conceptual boundary in a cloud-native package: a separate
process (or cluster) holds keys, enforces access policy, and exposes only
cryptographic operations through an authenticated HTTP API.  Vault's
\emph{transit secrets engine}~\cite{vault-transit} functions as a
``cryptography as a service'' layer with built-in key versioning, automatic
rotation, and a rich audit trail.

Integrating Vault into application-layer TLS or signing workflows is,
however, non-trivial.  Applications that use OpenSSL directly must either
call Vault's HTTP API themselves---adding a heavyweight dependency and a
significant refactoring burden---or depend on a thin wrapper library that
re-implements enough of the EVP\_PKEY interface to satisfy the caller.
Neither approach is satisfactory at scale.

We resolve this tension by exploiting the \emph{provider model} introduced
in OpenSSL~3.0~\cite{openssl3-providers}.  A provider is a dynamically
loaded module that implements one or more cryptographic operations through a
well-defined dispatch table.  Applications that use the standard
\texttt{EVP\_*} API are entirely unaware of the underlying provider; the
runtime selects the appropriate one based on algorithm properties and the
active library context.  By packaging the Vault integration as a provider,
we achieve full transparency: existing binaries, system TLS stacks, and
OpenSSL-based utilities obtain the isolation properties of Vault without any
code changes beyond a configuration file entry.

The contributions of this paper are:
\begin{itemize}
  \item A design and open-source implementation of \texttt{openssl-provider-vault},
        an OpenSSL~3 provider backed by the Vault transit engine
        (\S\ref{sec:design}, \S\ref{sec:impl}).
  \item An analysis of the security properties that the design provides, and
        the residual risks it does not eliminate (\S\ref{sec:security}).
  \item A discussion of the impedance mismatches between OpenSSL's streaming
        cryptography model and Vault's prehash-oriented REST API, and how we
        resolve them (\S\ref{sec:challenges}).
\end{itemize}

% ── 2. Background ────────────────────────────────────────────────────────────
\section{Background}
\label{sec:background}

\subsection{The OpenSSL~3 Provider Model}

OpenSSL~3 restructured its algorithm dispatch machinery around the concept of
\emph{providers}~\cite{openssl3-providers}.  A provider is a shared library
(or a statically compiled module) that exports a set of \emph{operation
implementations} through a C dispatch table.  Each operation corresponds to
an algorithm class---signature, asymmetric cipher, symmetric cipher, key
management, MAC, digest, and so on---and is identified by a combination of
algorithm name and property string.

The runtime selects a provider at algorithm instantiation time.  If an
application calls \texttt{EVP\_DigestSignInit} with algorithm
\texttt{"RSA-SHA256"} and the active library context has been configured with
\texttt{openssl-provider-vault}, the provider's \texttt{signature} dispatch table is
consulted first.  From the application's perspective, the call is
indistinguishable from the built-in software provider.

Key objects cross provider boundaries via a \emph{key reference} mechanism.
A provider exports keys as opaque byte sequences through
\texttt{keymgmt\_export} / \texttt{keymgmt\_import}; a receiving provider
reconstructs the local representation from those bytes.  This decoupling
lets openssl-provider-vault store only a lightweight handle to the remote key while
still participating in the standard \texttt{EVP\_PKEY} ownership model.

\subsection{HashiCorp Vault Transit Engine}

The Vault transit secrets engine~\cite{vault-transit} exposes cryptographic
operations over an authenticated HTTPS API.  Keys are created with a name
and a type (\texttt{rsa-2048}, \texttt{ecdsa-p256}, \texttt{ed25519},
\texttt{aes256-gcm96}, etc.) and are never exportable by default.  The
engine supports:

\begin{itemize}
  \item \textbf{Sign / Verify}: \texttt{POST /v1/transit/sign/:name} accepts a
        base64-encoded input (raw data or a digest), an optional hash algorithm,
        and a signature algorithm (\texttt{pss}, \texttt{pkcs1v15},
        \texttt{ecdsa}).
  \item \textbf{Encrypt / Decrypt}: symmetric and RSA operations;
        ciphertext is returned as a versioned envelope string
        (\texttt{vault:v\emph{N}:\emph{base64}}), where \emph{N} is the key
        version used for encryption.
  \item \textbf{HMAC}: \texttt{POST /v1/transit/hmac/:name/:alg}.
  \item \textbf{Key metadata}: \texttt{GET /v1/transit/keys/:name} returns
        the key type, version history, and PEM-encoded public keys.
\end{itemize}

Access is governed by Vault policies; every operation is recorded in the
audit log, giving operators a complete and tamper-evident record of key
usage.

% ── 3. Design ────────────────────────────────────────────────────────────────
\section{Design}
\label{sec:design}

\subsection{The Central Invariant: Keys Never Leave Vault}

The defining property of openssl-provider-vault is that \emph{no private key
material is ever present in the process memory of the application using
OpenSSL}.  The provider holds only a \emph{key reference} struct
(\texttt{vault\_keyref\_t}) containing:

\begin{itemize}
  \item the Vault key name and current version number,
  \item the key type string (\texttt{"rsa-2048"}, \texttt{"ecdsa-p256"},
        \textit{etc.}),
  \item optionally, a cached copy of the \emph{public} key in DER-encoded
        SubjectPublicKeyInfo format, fetched once at key creation or load time.
\end{itemize}

The public key is cached locally as a performance optimization for operations
that only require it (certificate verification, TLS handshake public-key
extraction) and carries no secret material.

When an operation requiring key material is requested---such as signing,
decrypting, or HMAC computation---the provider serializes the input according
to Vault's API contract, transmits it over an HTTPS connection with standard
TLS server authentication, and returns only the result (signature bytes,
plaintext, MAC tag) to the caller.  The private key is exercised exclusively
inside Vault.

\begin{figure}[t]
  \centering
  \begin{lstlisting}[language=C,caption={Key reference structure.},
                     label={lst:keyref}]
typedef struct {
    char          *key_name;
    int            key_version;
    char          *key_type;
    unsigned char *pubkey_der;
    size_t         pubkey_der_len;
} vault_keyref_t;
  \end{lstlisting}
\end{figure}

\subsection{Supported Operations}

openssl-provider-vault implements four provider operation classes:

\textbf{KEYMGMT.}  Handles key lifecycle: creation (\texttt{gen}), loading
from a serialized reference (\texttt{load}), parameter queries
(\texttt{get\_params}), and public-key export (\texttt{export}).  Private-key
export is unconditionally rejected; the provider returns an error if
\texttt{OSSL\_KEYMGMT\_SELECT\_PRIVATE\_KEY} is requested.  Key generation
maps to \texttt{POST /v1/transit/keys/:name}; the key name is supplied via a
custom \texttt{OSSL\_PARAM} (\texttt{"vault-key-name"}).

\textbf{SIGNATURE.}  Implements sign, verify, and the digest-sign/verify
streaming variants for RSA-PSS, RSA-PKCS\#1 v1.5, ECDSA, and Ed25519.

\textbf{ASYM\_CIPHER.}  Wraps Vault's \texttt{encrypt} and \texttt{decrypt}
endpoints for RSA-based asymmetric encryption.

\textbf{MAC.}  Delegates HMAC computation to Vault's \texttt{hmac} endpoint.

\textbf{STORE.}  Resolves \texttt{vault:} URIs to key references.
Applications load a Vault-backed key through OpenSSL's STORE framework using
a URI such as \texttt{"vault:my-signing-key"}, analogous to other
URI-backed loaders such as \texttt{pkcs11:}.

\subsection{Configuration}
\label{sec:config}

The provider is activated by adding a stanza to \texttt{openssl.cnf}:

\begin{lstlisting}[caption={Minimal openssl.cnf for openssl-provider-vault.}]
[openssl_init]
providers = provider_sect

[provider_sect]
vault = vault_sect

[vault_sect]
module = /usr/lib/ossl-modules/vault.so
vault_addr  = https://vault.example.com:8200
vault_token = s.xxxxx
\end{lstlisting}

No application code changes are required beyond configuring OpenSSL to load
the provider and supply the necessary Vault connection parameters.

% ── 4. Implementation ────────────────────────────────────────────────────────
\section{Implementation}
\label{sec:impl}

openssl-provider-vault is implemented in approximately 2\,600 lines of C99 across
ten source modules (Table~\ref{tab:modules}).  It depends on libcurl for
HTTP transport and json-c for response parsing; all other dependencies are
satisfied by the OpenSSL library itself.

\begin{table}[t]
  \caption{Source modules and their roles.}
  \label{tab:modules}
  \small
  \begin{tabularx}{\columnwidth}{lX}
    \toprule
    Module & Responsibility \\
    \midrule
    \texttt{provider.c}          & Provider entry point, context lifecycle \\
    \texttt{vault\_client.c}     & HTTP client (libcurl), JSON parsing \\
    \texttt{vault\_format.c}     & Base64 encode/decode, Vault wire format \\
    \texttt{vault\_keyref.c}     & Key reference struct, serialization \\
    \texttt{vault\_store\_uri.c} & \texttt{vault:} URI parser \\
    \texttt{keymgmt.c}           & KEYMGMT dispatch implementation \\
    \texttt{signature.c}         & SIGNATURE dispatch implementation \\
    \texttt{asym\_cipher.c}      & ASYM\_CIPHER dispatch implementation \\
    \texttt{mac.c}               & MAC (HMAC) dispatch implementation \\
    \texttt{store.c}             & STORE dispatch implementation \\
    \bottomrule
  \end{tabularx}
\end{table}

\subsection{Testing}

At the time of writing, the test suite comprises roughly 3\,500 lines of
test code and just over 150 unit and integration tests across nine
binaries. It is designed so that every test binary is entirely self-contained:
each links the source files it exercises directly, rather than loading the
provider shared library.  Where a test exercises a module that normally calls
vault\_client.c, a \emph{mock vault client} is substituted, allowing
full control over Vault API responses via cmocka's \texttt{will\_return}
mechanism~\cite{cmocka}.

One test binary (\texttt{test\_vault\_client}) is an integration test that
starts a local Vault dev server automatically if the \texttt{vault(1)} binary
is present on \texttt{PATH}, or connects to an external server if
\texttt{VAULT\_ADDR}/\texttt{VAULT\_TOKEN} are set.  If neither condition is
met, the binary exits with code~77 (autotest ``skip'').

\textbf{OOM / allocation-failure testing.}  Heap exhaustion is a
common cause of subtle memory-safety bugs in C code.  To exercise all
allocation failure paths, we use GNU ld's \texttt{--wrap} mechanism to
intercept \texttt{malloc} and \texttt{calloc} at link time.  A shared helper
(\texttt{wrap\_malloc.c}) exposes a \texttt{wrap\_malloc\_fail\_after(n)}
function that causes the next allocation from compiled-in code to return
\texttt{NULL} after $n$ successful calls.  Because the system libraries
(libcrypto, libcurl, libc) are shared objects, their internal allocations
are \emph{not} intercepted---only direct calls from provider source files
are affected, giving precise control over which allocation is under test.

% ── 5. Design Challenges ─────────────────────────────────────────────────────
\section{Design Challenges}
\label{sec:challenges}

\subsection{The Prehashing Impedance Mismatch}

OpenSSL's digest-sign API is a streaming interface: the caller feeds data
incrementally through \texttt{Digest\-Sign\-Update} and finalizes with
\texttt{Digest\-Sign\-Final}.  Vault's sign endpoint, by contrast, accepts
a single pre-digested input alongside the \texttt{prehashed=true} flag.

The provider resolves this by maintaining an \texttt{EVP\_MD\_CTX} inside
the signature context struct.  During update calls, data is accumulated into
the local digest context; at finalization time, the digest is computed
locally and the resulting hash is submitted to Vault.  This is correct and
secure: the hash is computed in the same process that would otherwise compute
it before handing off to any software signer; no trust boundary is crossed.

\subsection{Context Duplication (\texttt{dupctx})}

The OpenSSL provider interface requires that a context can be duplicated at
any point during a streaming operation via \texttt{dupctx}.  For the
signature provider, this means copying the \texttt{EVP\_MD\_CTX} mid-stream.
The implementation calls \texttt{EVP\_MD\_CTX\_copy\_ex} on the source
context and attaches the result to the newly allocated destination context.
Care must be taken to null-initialize the destination's \texttt{mdctx}
pointer before branching on the source, so that a copy-failure path calls
\texttt{EVP\_MD\_CTX\_free(NULL)} (a no-op) rather than double-freeing a
partially initialized struct.

\subsection{Vault Wire Format}

Vault's transit engine encodes all binary data in URL-safe base64 (no
padding in some contexts) and prefixes ciphertext with a versioned envelope
(\texttt{vault:v\emph{N}:\emph{base64}}), where \emph{N} is the key
version number.  The \texttt{vault\_format.c} module
centralizes all encoding and decoding using OpenSSL's BIO chain API, ensuring
a single, tested code path for all format conversions.

\subsection{Algorithm Identifier Generation}

Some callers (notably those constructing X.509 certificate signing requests)
query the signature context for an ASN.1 \texttt{AlgorithmIdentifier} via
the \texttt{OSSL\_SIGNATURE\_PARAM\_ALGORITHM\_ID} parameter.  The provider
synthesizes the correct DER-encoded structure using \texttt{X509\_ALGOR\_new}
and \texttt{i2d\_X509\_ALGOR}, driven by the hash NID and key type stored in
the context.  This is necessary for compatibility with OpenSSL's CSR and
certificate-generation code paths, which retrieve the algorithm identifier
from the signer context rather than from the key.

% ── 6. Security Analysis ─────────────────────────────────────────────────────
\section{Security Properties and Limitations}
\label{sec:security}

\subsection{What openssl-provider-vault Protects}

\textbf{Memory-resident key material.}  Because the private key is never
retrieved from Vault, a process-level compromise---including a full heap or
stack dump, an arbitrary-read vulnerability, or a debugger attach---cannot
yield the private key.  This is the primary threat addressed.

\textbf{Key-use audit trail.}  Every signing, encryption, and HMAC operation
generates an entry in Vault's audit log.  Operators can detect unexpected
key usage, enforce rate limits via Vault policies, and obtain a forensic
timeline after an incident.

\textbf{Key revocation and rotation.}  Vault supports immediate key
revocation and seamless key rotation without re-keying existing ciphertexts
(for symmetric keys).  The \texttt{vault\_keyref\_t} stores the key version
used at operation time, enabling Vault's min-decryption-version policy to
enforce forward secrecy windows.

\textbf{Access control.}  Vault policies can restrict which application
identities (AppRole, Kubernetes auth, etc.) may use which keys for which
operations, with much finer granularity than file system permissions.

\subsection{Residual Risks}

\textbf{Vault availability.}  openssl-provider-vault introduces a synchronous
dependency on the Vault cluster for every private-key operation.  A Vault
outage or network partition will prevent signing and decryption.  Operators
should deploy Vault with its recommended HA configuration~\cite{vault-ha}.

\textbf{Transport security.}  The provider communicates with Vault over
HTTPS.  The security of this channel depends on correct certificate
validation and the absence of a MITM adversary.  Currently the provider
relies on libcurl's default server-certificate validation; future work
should expose configuration for a pinned CA bundle and, optionally,
client-certificate-based authentication.

\textbf{Token storage.}  The Vault token is currently read from environment
variables or \texttt{openssl.cnf}.  An adversary with read access to either
can obtain a token and use it to request signatures directly.  More
robust authentication methods (e.g., Vault Agent, Kubernetes auth, AWS IAM
auth) should be preferred in production, and token storage should be treated
with the same care as the key it protects.

\textbf{Side channels in Vault.}  openssl-provider-vault does not address timing
or other side channels within the Vault process itself.  Vault uses
Go's standard library crypto, which incorporates timing countermeasures---RSA
blinding, constant-time comparison via \texttt{crypto/subtle}---for critical
operations.  The extent to which Go's big-number arithmetic provides complete
constant-time guarantees remains an area of active research~\cite{go-constant-time-bignum};
this dimension of side-channel resistance is outside the scope of the present work.

\textbf{Scope: no FIPS certification.}  openssl-provider-vault is not a FIPS
140-validated module.  Organizations with FIPS requirements should evaluate
whether Vault's own FIPS variant satisfies their compliance needs.

% ── 7. Related Work ──────────────────────────────────────────────────────────
\section{Related Work}
\label{sec:related}

\textbf{PKCS\#11-based providers.}
\texttt{pkcs11-provider}~\cite{pkcs11-provider} is the primary prior art for
OpenSSL~3 provider-based key isolation.  It implements the same drop-in
transparency goal against PKCS\#11 tokens---physical HSMs, SoftHSM, TPMs.
openssl-provider-vault is complementary: it targets deployments where a cloud-native
secret-management system is preferred over a hardware device, and where
Vault's policy engine and audit capabilities provide value beyond raw
key isolation.

\textbf{Post-quantum providers.}
The Open Quantum Safe project's \texttt{oqs-provider}~\cite{oqs-provider}
demonstrates the same OpenSSL~3 provider dispatch pattern in a different
domain: it registers quantum-safe KEM and signature algorithms so that
TLS~1.3, X.509, and S/MIME consumers can use them transparently.  While
its focus is algorithmic diversity rather than key isolation, it is a
notable example of the provider model's generality and of the ecosystem
of third-party providers that have emerged since OpenSSL~3.0.

\textbf{Engine-based approaches.}  OpenSSL's deprecated engine interface was
used by numerous projects (e.g., \texttt{engine\_pkcs11},
\texttt{aws-lc}) to redirect private-key operations.  The provider model is
the supported successor and avoids the thread-safety and context-propagation
problems inherent in the global engine registry.

\textbf{Vault PKI / TLS integration.}  HashiCorp provides several
Vault-adjacent tools for TLS (Vault Agent, Envoy SDS, cert-manager) but
these operate at the certificate issuance layer, not at the signing primitive
layer.  They do not prevent key material from residing in the application
process after issuance.  openssl-provider-vault addresses the residual risk of
an already-issued key being exposed at runtime.

% ── 8. Future Work ───────────────────────────────────────────────────────────
\section{Future Work}
\label{sec:future}

\textbf{Performance characterization.}  A systematic latency and throughput
study across key types, Vault deployment topologies (local, Vault Cloud, HA
cluster), and TLS session reuse strategies is needed to quantify the
performance overhead of the remote-dispatch model.

\textbf{mTLS and certificate-based Vault authentication.}  Replacing token
authentication with a client-certificate-based Vault auth method would
eliminate the token-storage risk identified in \S\ref{sec:security}.

\textbf{Key derivation context.}  Vault's \texttt{derived: true} key flag
enables per-context key derivation.  Exposing this through a custom
\texttt{OSSL\_PARAM} would support authenticated encryption schemes where
each session derives a unique key from a shared root.

\textbf{Asynchronous signing.}  Long-lived TLS servers that sign thousands
of handshakes per second must serialize all signing operations through the
Vault HTTP endpoint.  Exposing an asynchronous signing interface and
pipelining requests over a persistent connection pool would substantially
reduce latency under load.

\textbf{Digest provider.}  While hashing inside Vault adds an unnecessary
network round-trip for most use cases, the \texttt{transit/hash} endpoint
could be exposed as an optional provider-backed DIGEST operation for
regulated environments that require all cryptographic operations to be
recorded in the Vault audit log.

\textbf{PKCS\#11 compatibility layer.}  An open-source Vault-backed
PKCS\#11 module would be a natural complement to the provider design for
software ecosystems that expect HSM-style interfaces, particularly for
deployments that wish to avoid proprietary middleware in their cryptographic
stack.  Such a module could expose Vault-managed keys through the de facto
standard token API used by existing HSM integrations, while reusing much of
the same remote-operation model described in this work.

% ── 9. Conclusion ────────────────────────────────────────────────────────────
\section{Conclusion}
\label{sec:conclusion}

We have presented openssl-provider-vault, an OpenSSL~3 provider that delegates
private-key operations to HashiCorp Vault's transit secrets engine.  By
implementing the standard OpenSSL provider dispatch interface, the system
achieves drop-in transparency: any application or library built on the
OpenSSL EVP API---from \texttt{openssl(1)} to production TLS servers---gains
the isolation and auditability of Vault-managed keys without modification.
The provider handles the prehashing impedance mismatch between OpenSSL's
streaming interface and Vault's prehash-oriented REST API, correctly manages
context duplication mid-stream, and generates standard ASN.1 algorithm
identifiers required by certificate-issuance code paths.

The design ensures that private key material is never present in application
process memory; compromise of the application, its heap, or its TLS library
cannot yield the key.  We have described the residual risks---chiefly Vault
availability, token storage, and transport security---and outlined mitigations
for each.

openssl-provider-vault is open-source and available at
\url{https://github.com/sorenisanerd/openssl-provider-vault}.

% ── References ───────────────────────────────────────────────────────────────
\bibliographystyle{ACM-Reference-Format}
\begin{thebibliography}{9}

\bibitem{vault-transit}
HashiCorp.
\newblock Vault Transit Secrets Engine Documentation.
\newblock \url{https://developer.hashicorp.com/vault/docs/secrets/transit},
  2024.

\bibitem{openssl3-providers}
OpenSSL Project.
\newblock Providers, OpenSSL 3.x Manual.
\newblock \url{https://www.openssl.org/docs/man3.0/man7/provider.html}, 2024.

\bibitem{pkcs11-provider}
Simo Sorce et al.
\newblock pkcs11-provider: A PKCS\#11 Provider for OpenSSL~3.
\newblock \url{https://github.com/openssl-projects/pkcs11-provider}, 2023.

\bibitem{cmocka}
Andreas Schneider and Google Inc.
\newblock cmocka: Unit testing framework for C.
\newblock \url{https://cmocka.org}, 2013.

\bibitem{oqs-provider}
Open Quantum Safe Project.
\newblock oqs-provider: An OpenSSL 3 Provider for Quantum-Safe Cryptography.
\newblock \url{https://github.com/open-quantum-safe/oqs-provider}, 2025.

\bibitem{go-constant-time-bignum}
Lucas Meier.
\newblock Constant-Time Big Numbers (for Go).
\newblock EPFL Technical Report, 2021.
\newblock \url{https://www.epfl.ch/labs/dedis/wp-content/uploads/2021/07/report-2021-1-LucasMeier_ConstantTime.pdf}

\bibitem{vault-ha}
HashiCorp.
\newblock Vault High Availability Documentation.
\newblock \url{https://developer.hashicorp.com/vault/docs/concepts/ha}, 2024.

\end{thebibliography}

\end{document}
